% GENERATED FILE. Edit canonical lessons, then run npm run generate:books.
% canonical-source-hash: fnv1a64:ed1e32c303d0d455
% canonical-lessons: FR-C403-pompe, FR-C403-gazole, FR-C403-carburant, FR-C403-reservoir, FR-C403-demarrer

\chapter{À la pompe}
\label{ch:fr-a-la-pompe}
\hlchaptermodality{403}

\begin{chapteropening}
\textbf{By the end of this chapter:} \emph{I can fill up a car: the pump, diesel, fuel, the tank and starting off again.}

Complete the last lesson of chapter 403: I can fill up a car: the pump, diesel, fuel, the tank and starting off again.
\end{chapteropening}

\section[la pompe]{la pompe — the pump}

\label{lesson:FR-C403-pompe}



\begin{quote}
Before the new one: say the French for to honk, then the French for to mist up.
\end{quote}

\subsection*{You'll want to know: la pompe}
\textbf{la pompe} — \textquotedblleft{}the pump\textquotedblright{}.

\begin{quote}
La pompe numéro trois ne marche pas.
\end{quote}

\begin{quote}
\emph{Pump number three is not working.}
\end{quote}

\emph{La pompe à essence} is the petrol pump; \emph{une pompe à vélo} is a bicycle pump.

\begin{cousinweb}
The origin of \emph{pompe} is not settled. It is usually linked to Middle Low German \emph{pumpe} 'pump' (EtymDB) and may imitate the sound of the pump.
\end{cousinweb}

\begin{grammarlens}[title={What you've built}]
One more word for telling, explaining and describing at B1.
\end{grammarlens}

\subsection*{Your turn}
\begin{itemize}
\raggedright
  \item \emph{Say it:} \emph{la pompe}
  \item \emph{Say it:} \emph{la pompe}, once more
  \item \emph{Say it:} a sentence of your own with \emph{la pompe}
  \item \emph{Recall:} say the French for to honk, then the French for to mist up, then say \emph{la pompe} again
\end{itemize}

\begin{tcolorbox}[breakable,colback=teal!4,colframe=teal!35!black,title={Before you move on}]
What does la pompe mean? (\textquotedblleft{}The pump\textquotedblright{}.) Say it once more.
\end{tcolorbox}
\section[le gazole]{le gazole — the diesel (fuel)}

\label{lesson:FR-C403-gazole}



\begin{quote}
Before the new one: say the French for to mist up, then the French for the pump.
\end{quote}

\subsection*{You'll want to know: le gazole}
\textbf{le gazole} — \textquotedblleft{}the diesel (fuel)\textquotedblright{}.

\begin{quote}
Ma voiture roule au gazole.
\end{quote}

\begin{quote}
\emph{My car runs on diesel.}
\end{quote}

Also spelled \emph{le gasoil}; many people say \emph{le diesel}.

\begin{cousinweb}
\emph{gazole} is usually explained as a French spelling of English \emph{gas oil}; EtymDB does not list it.
\end{cousinweb}

\begin{grammarlens}[title={What you've built}]
One more word for telling, explaining and describing at B1.
\end{grammarlens}

\subsection*{Your turn}
\begin{itemize}
\raggedright
  \item \emph{Say it:} \emph{le gazole}
  \item \emph{Say it:} \emph{le gazole}, once more
  \item \emph{Say it:} a sentence of your own with \emph{le gazole}
  \item \emph{Recall:} say the French for to mist up, then the French for the pump, then say \emph{le gazole} again
\end{itemize}

\begin{tcolorbox}[breakable,colback=teal!4,colframe=teal!35!black,title={Before you move on}]
What does le gazole mean? (\textquotedblleft{}The diesel (fuel)\textquotedblright{}.) Say it once more.
\end{tcolorbox}
\section[le carburant]{le carburant — the fuel}

\label{lesson:FR-C403-carburant}



\begin{quote}
Before the new one: say the French for the pump, then the French for the diesel.
\end{quote}

\subsection*{You'll want to know: le carburant}
\textbf{le carburant} — \textquotedblleft{}the fuel\textquotedblright{}.

\begin{quote}
Le prix du carburant augmente encore.
\end{quote}

\begin{quote}
\emph{The price of fuel is going up again.}
\end{quote}

The general word: petrol and diesel are both \emph{carburant}.

\begin{cousinweb}
\emph{carburant} is built on \emph{carburer} 'to mix fuel with air', from \emph{carbure} 'carbide', a learned formation on Latin \emph{carbō} 'charcoal' (Lewis \& Short).
\end{cousinweb}

\begin{grammarlens}[title={What you've built}]
One more word for telling, explaining and describing at B1.
\end{grammarlens}

\subsection*{Your turn}
\begin{itemize}
\raggedright
  \item \emph{Say it:} \emph{le carburant}
  \item \emph{Say it:} \emph{le carburant}, once more
  \item \emph{Say it:} a sentence of your own with \emph{le carburant}
  \item \emph{Recall:} say the French for the pump, then the French for the diesel, then say \emph{le carburant} again
\end{itemize}

\begin{tcolorbox}[breakable,colback=teal!4,colframe=teal!35!black,title={Before you move on}]
What does le carburant mean? (\textquotedblleft{}The fuel\textquotedblright{}.) Say it once more.
\end{tcolorbox}
\section[le réservoir]{le réservoir — the tank (for fuel); the reservoir}

\label{lesson:FR-C403-reservoir}



\begin{quote}
Before the new one: say the French for the diesel, then the French for the fuel.
\end{quote}

\subsection*{You'll want to know: le réservoir}
\textbf{le réservoir} — \textquotedblleft{}the tank (for fuel); the reservoir\textquotedblright{}.

\begin{quote}
Le réservoir est presque vide.
\end{quote}

\begin{quote}
\emph{The tank is almost empty.}
\end{quote}

\emph{Faire le plein} — to fill the tank.

\begin{cousinweb}
\emph{réservoir} is built on \emph{réserver} 'to keep aside', a learned borrowing from Latin \emph{reservāre}, from \emph{servāre} 'to keep safe' (Lewis \& Short).
\end{cousinweb}

\begin{grammarlens}[title={What you've built}]
One more word for telling, explaining and describing at B1.
\end{grammarlens}

\subsection*{Your turn}
\begin{itemize}
\raggedright
  \item \emph{Say it:} \emph{le réservoir}
  \item \emph{Say it:} \emph{le réservoir}, once more
  \item \emph{Say it:} a sentence of your own with \emph{le réservoir}
  \item \emph{Recall:} say the French for the diesel, then the French for the fuel, then say \emph{le réservoir} again
\end{itemize}

\begin{tcolorbox}[breakable,colback=teal!4,colframe=teal!35!black,title={Before you move on}]
What does le réservoir mean? (\textquotedblleft{}The tank (for fuel); the reservoir\textquotedblright{}.) Say it once more.
\end{tcolorbox}
\section[démarrer]{démarrer — to start up (a car, an engine); to get going}

\label{lesson:FR-C403-demarrer}



\begin{quote}
Before the new one: say the French for the fuel, then the French for the tank.
\end{quote}

\subsection*{You'll want to know: démarrer}
\textbf{démarrer} — \textquotedblleft{}to start up (a car, an engine); to get going\textquotedblright{}.

\begin{quote}
Ce matin, la vieille voiture ne démarre pas.
\end{quote}

\begin{quote}
\emph{This morning the old car will not start.}
\end{quote}

Also for projects: \emph{le travail démarre lundi}.

\begin{cousinweb}
\emph{démarrer} is built from \emph{dé-} and \emph{amarrer} 'to moor', which is usually traced to Middle Dutch \emph{aenmarren} 'to tie up': to start was first to untie a boat. EtymDB lists the verb without a source.
\end{cousinweb}

\begin{grammarlens}[title={What you've built}]
One more word for telling, explaining and describing at B1.
\end{grammarlens}

\subsection*{Your turn}
\begin{itemize}
\raggedright
  \item \emph{Say it:} \emph{démarrer}
  \item \emph{Say it:} \emph{démarrer}, once more
  \item \emph{Say it:} a sentence of your own with \emph{démarrer}
  \item \emph{Recall:} say the French for the fuel, then the French for the tank, then say \emph{démarrer} again
\end{itemize}

\begin{tcolorbox}[breakable,colback=teal!4,colframe=teal!35!black,title={Before you move on}]
What does démarrer mean? (\textquotedblleft{}To start up (a car, an engine); to get going\textquotedblright{}.) Say it once more.
\end{tcolorbox}
